% Generated draft: requires mathematical and source review.
\documentclass[11pt]{article}
\usepackage[margin=25mm]{geometry}
\usepackage{fontspec}
\setmainfont{Latin Modern Roman}
\usepackage{amsmath,amssymb,mathtools,unicode-math}
\setmathfont{Latin Modern Math}
\usepackage{enumitem,needspace,xurl,hyperref}
\setlength{\parindent}{0pt}
\setlength{\parskip}{5pt}
\setlength{\emergencystretch}{3em}
\newcommand{\N}{\mathbb{N}}
\newcommand{\Z}{\mathbb{Z}}
\newcommand{\Q}{\mathbb{Q}}
\newcommand{\R}{\mathbb{R}}
\newcommand{\C}{\mathbb{C}}
\newcommand{\F}{\mathbb{F}}
\newcommand{\A}{\mathbb{A}}
\newcommand{\PP}{\mathbb{P}}
\newcommand{\E}{\mathbb{E}}
\newcommand{\eps}{\varepsilon}
\newcommand{\dd}{\,\mathrm{d}}
\begin{document}
\section*{UnsolvedMath ID: 20002075}
\textbf{Source problem code:} catalogue entry
\subsection{Definitions and mathematical statement}
Let \(\omega\) be an integral symplectic form on \(T^4\), with the integral normalization understood so that \([\omega]\in H^2(T^4;\Z)\). Let \(\pi:P^5\to T^4\) be the principal circle bundle with \(c_1(P)=[\omega]\), equipped with its prequantization contact structure \(\xi=\ker\alpha\), where the connection curvature is the pullback of the appropriately normalized \(\omega\). The source's Question \(1.17\) conjectures that \((P,\xi)\) has no Stein filling: no compact complex three-dimensional Stein domain has this contact boundary.

The aid reports a stronger topological statement: the total space of every principal circle bundle over \(T^4\) admits no six-dimensional Stein or Weinstein filling, regardless of contact structure on that total space. Its stated obstruction uses that \(P\) is a closed oriented aspherical five-manifold, so \(P\simeq B\pi_1(P)\) and its classifying map sends \([P]\) nontrivially to group homology. A six-dimensional Stein/Weinstein domain has handles of index at most three; viewed relative to its entire boundary \(P\), this implies \(\pi_1(P)\to\pi_1(W)\) is an isomorphism. A classifying map then extends across \(W\), forcing the boundary fundamental class to map to zero, a contradiction. The filling has boundary exactly \(P\), with the usual orientations. This is the aid's reported resolution, not independent mathematical review of its topology inputs.

Its additional integral-cohomology proposition concerns the symplectic Euler class \(e=du\), where \(d\geq1\), \(u\in H^2(T^4;\Z)\) is primitive, and
\[
 Q=\langle u\smile u,[T^4]\rangle\ne0.
\]
For \(P_e\) define \(a=\pi^*u\in H^2(P_e;\Z)\). It generates the torsion summand \(\Z/d\Z\), trivial when \(d=1\), and the aid asserts
\[
 a^2\ne0\quad\Longleftrightarrow\quad d\nmid Q.
\]
In the integral alternating-form normal form, choosing an oriented basis \(x_1,\ldots,x_4\) of \(H^1(T^4;\Z)\),
\[
 u=x_1\smile x_2+m\,x_3\smile x_4,\qquad Q=2m,
\]
so the test is \(d\nmid2m\). This cup-square refinement is distinct from the fundamental-class nonfillability obstruction, which does not require that square to be nonzero.
\subsection{Short English statement}
The aid reports that torus circle-bundle total spaces are aspherical and cannot bound six-dimensional Stein or Weinstein domains, so the prequantization conjecture follows. It also gives a separate divisibility test for the square of a torsion cohomology class; nonfillability does not depend on that square being nonzero.
\subsection{Sources}
\begin{itemize}
\item Local UnsolvedMath registry, permanent ID \(20002075\), source code \texttt{catalogue entry}; canonical collection \texttt{aim-geometry-notes.json}, notes[412], zero-based index \(412\).
\item AIM workshop ``Holomorphic curves in contact geometry,'' Question \(1.17\): \url{https://aimath.org/WWN/contactgeom2/contactgeom2.pdf}. The source specifies a contact circle bundle over \(T^4\) with first Chern class the symplectic class. No PDF was opened in this session.
\item Supplied research attempt \(1\), result \texttt{full\_solution}, run status ``solved,'' rigor label ``proved,'' novelty confidence ``low.'' Registry status ``solved,'' classification \texttt{SOLVED-IN-LITERATURE}, updated 14 August 2026. Its candidate contribution is the torsion cup-square proposition in addition to the reported topological resolution.
\item No full authoritative references for the Stein handle theorem, classifying-space obstruction or integral normal-form calculation are supplied. Original remarks and literature fields are empty, with no proposer or year.
\end{itemize}
\end{document}
